\begin{lemma}\label{lem:nonic83-pole-degrees-direct}
The modular units $x,y$ used in the order-$83$ canonical recovery
have pole degrees $90,95$, respectively, on $X_1(83)_{\mathbb F_2}$.
These degrees can be certified independently of a plane equation
and of the canonical unit relations.
\end{lemma}
\begin{proof}
The unique Tate normal form and the birational changes
\[
 r=b/c,\quad s=c^2/(b-c),\quad
 x=(s-r)/(rs-2r+1),\quad y=(rs-2r+1)/(s^2-s-r+1)
\]
are given in \cite[Sections~2 and~2.1]{DvH}.
Their rational inverses are checked identically over $\mathbb Q$.
In the notation of that paper, $x=f_7/f_8$ and $y=f_8/f_9$.
These are modular units at level $83$: the vanishing equations for
orders $7,8,9$ cannot hold for a smooth point of exact order $83$.
The same argument applies in characteristic $2$, since the
prime-to-$2$ level point still has exact order $83$.
Thus only cusp valuations enter their pole degrees.

There are $41$ geometric cusps of width $83$, represented by
$E(q^{83})$ with Tate point $u=q^a$, $1\leq a\leq41$.
Exact expansions through precision $1200$, with explicit retained
precision after every rational operation, give
\[
 \sum_{a=1}^{41}\max\{0,-v_a(x)\}=90,\qquad
 \sum_{a=1}^{41}\max\{0,-v_a(y)\}=95.
\]
All $41$ pairs of orders are retained and replayed.
The elliptic equation is checked at every Tate point before the
normal-form changes are applied.

The other $41$ geometric cusps have width $1$ and are represented
by $E(q)$ with $u=\zeta_{83}$, modulo inversion.
Since $2$ has order $82$ modulo $83$, they form one Frobenius orbit
modulo inversion. Work exactly in
$\mathbb F_{2^{82}}=\mathbb F_2[z]/(1+z+\cdots+z^{82})$.
At $q=0$ the Tate point has coordinates
\[
 x_P=\frac{z}{(1+z)^2},\qquad
 y_P=\frac{z^2}{(1+z)^3}.
\]
Substituting these into the normal-form formulas gives finite,
nonzero values for both units. These values are fixed by
$z\mapsto z^{2^{41}}$, as required by inversion.
Every width-$1$ valuation is therefore zero.
This proves the asserted pole degrees without using a plane model
or a canonical relation.
\end{proof}
